\chapter{Step 1: reports into labels}
\label{ch:labels}
\whereami{2}

\begin{tldr}
\begin{itemize}
\item We used an open \textbf{large language model} (LLM), running inside Kaggle, to read all 4,407 reports.
\item \textbf{v3} (Qwen2.5-7B) asked for one word per finding and used ``any grade counts'': agreement with the
      doctors' gold labels \textbf{0.809} AUC.
\item \textbf{v4} (Qwen2.5-32B) used the organisers' thresholds, five answer letters, and read the model's
      \textbf{probabilities} instead of its text: \textbf{0.879}. Leaderboard of the image model: 0.889 $\rightarrow$ 0.905.
\item Reports cannot label what they do not mention: \textbf{Synovitis} is silent in 84\,\% of them.
\end{itemize}
\end{tldr}

\section{What a language model does}

A language model reads text split into \term{tokens} (word pieces) and, at every position, outputs a probability
for every possible next token. Generating text means picking a token, appending it and repeating.

\begin{figure}[H]
\centering
\begin{tikzpicture}[x=1mm,y=1mm]
  \node[draw=grid, fill=white, rounded corners=1.5mm, text width=56mm, font=\scriptsize\ttfamily, align=left,
        inner sep=2mm] (p) at (0,0) {[instructions: 12 findings, definitions, letters P M U A N]\\[1mm]
        [report of Study K]\\[1mm] \{"ACL": "};
  \node[draw=accent, fill=keyfill, rounded corners=2mm, minimum width=20mm, minimum height=12mm, font=\footnotesize,
        align=center] (m) at (50,0) {\faRobot\\ LLM};
  \draw[-{Stealth}, muted, line width=0.8pt] (p) -- (m);
  \foreach \t/\v [count=\i from 0] in {P/0.95, M/0.02, U/0.02, A/0.01, N/0.00, {the}/0.00} {
    \node[font=\scriptsize\ttfamily, anchor=east] at (72,{9-\i*3.6}) {\t};
    \fill[grid] (73,{9-\i*3.6-1.1}) rectangle ++(40,2.2);
    \fill[us] (73,{9-\i*3.6-1.1}) rectangle ++({40*\v},2.2);
    \node[font=\scriptsize, anchor=west] at (114,{9-\i*3.6}) {\v};
  }
  \draw[-{Stealth}, muted, line width=0.8pt] (m) -- (66,0);
  \node[font=\scriptsize, text=muted, align=center] at (93,-15) {probability of each possible next token\\ (the full vocabulary has \ca 150,000)};
\end{tikzpicture}
\caption{The prompt ends right where the answer for ACL should start. The model's next-token probabilities
\emph{are} the answer: 95\,\% ``P'' (present, past the threshold).}
\label{fig:nexttoken}
\end{figure}

\begin{wordcheck}
\term{LLM}: large language model, e.g.\ ChatGPT, Claude, Qwen. \term{Open-weights model}: one whose weights you can
download and run yourself (Qwen is made by Alibaba). \term{Prompt}: the text you give the model.
\term{4-bit quantization}: storing each weight in 4 bits instead of 16 so a big model fits in less memory.
\end{wordcheck}

\section{Why an open model, inside Kaggle?}

We believed the competition's data rules forbade sending report text to hosted services such as ChatGPT or
Claude, so the reports never left Kaggle. A Kaggle notebook offers two T4 GPUs with 16\,GB each, for at most
12 hours per session. A 32-billion-parameter model only fits in 4-bit form (19\,GB spread over both GPUs).

\begin{watchout}
That belief was \textbf{wrong}. The organisers had stated in a pinned forum post that hosted LLM APIs, including
for labelling reports, were allowed. We never read that post. The winners used the strongest hosted models and,
crucially, let them look at the images too (Chapter~\ref{ch:winlabels}). Lesson: \textbf{check the rules threads on
day one, and treat every self-imposed rule as a question to verify.}
\end{watchout}

\section{Version 3: the first labeller}

\begin{center}\small
\begin{tabularx}{\linewidth}{l>{\raggedright\arraybackslash}X}
\toprule
Model & Qwen2.5-7B-Instruct (16-bit) \\
Prompt & ``for each finding, answer \code{present}, \code{minimal}, \code{uncertain}, \code{absent} or \code{not\_mentioned}; any grade counts'' \\
Reading the answer & parse the generated text; map to label values 1, 0.5 or 0 \\
Runtime & 3.8 hours for 4,407 reports; 3.6\,\% of reports had an unparseable answer \\
Agreement with gold & \textbf{macro AUC 0.809} \\
\bottomrule
\end{tabularx}
\end{center}

Two problems showed up when we compared it with the 58 gold studies:
\begin{enumerate}
\item \textbf{Wrong definitions.} ``Any grade counts'' turned every small effusion and every mild sprain into a
      positive. The organisers' rule is the opposite (Figure~\ref{fig:thresholds}).
\item \textbf{Silence became ``absent''.} The model almost never answered \code{not\_mentioned} (0--2 of 58 gold
      reports per finding). But ``the report does not mention synovitis'' is not the same as ``no synovitis''.
\end{enumerate}

\section{Version 4: four changes}

\subsection*{1. The organisers' thresholds, written into the prompt}

For every finding the prompt lists what counts as \textbf{P} (present and past the threshold) and what counts as
\textbf{M} (present but below it). Two examples:

\begin{center}\small
\begin{tabularx}{\linewidth}{l>{\raggedright\arraybackslash}X>{\raggedright\arraybackslash}X}
\toprule
 & P (counts) & M (does not count) \\
\midrule
Effusion & moderate or large joint effusion & small, trace, minimal, mild or physiological fluid \\
MCL & acute high-grade partial or complete tear (grade 2--3) & grade 1 sprain, edema around an intact ligament, old scarring \\
\bottomrule
\end{tabularx}
\end{center}

\subsection*{2. Five answers instead of three}

\begin{figure}[H]
\centering
\begin{tikzpicture}[x=1mm,y=1mm,
  let/.style={draw=#1, fill=#1!15, rounded corners=1.5mm, minimum width=24mm, minimum height=19mm, align=center,
              font=\footnotesize, text width=22mm}]
  \node[let=them] at (0,0) {\textbf{\Large P}\\ present, past the threshold};
  \colorlet{mcol}{them!60}
  \node[let=mcol] at (27,0) {\textbf{\Large M}\\ present, below the threshold};
  \node[let=whyline] at (54,0) {\textbf{\Large U}\\ unclear: no size, or hedged};
  \node[let=us] at (81,0) {\textbf{\Large A}\\ explicitly absent / normal};
  \node[let=muted] at (108,0) {\textbf{\Large N}\\ not mentioned at all};
  \foreach \x/\v in {0/1, 27/{\ca 0.2}, 54/{\ca 0.5}, 81/0, 108/{measured}} {
    \node[font=\scriptsize, text=ink] at (\x,-14) {value: \textbf{\v}};
  }
\end{tikzpicture}
\caption{The five answer letters and their label values. The values for M, U and N were measured on the gold
studies per finding (how often a gold doctor said ``positive'' after that answer), pulled towards the defaults
shown here so that 58 studies do not overrule common sense.}
\label{fig:letters}
\end{figure}

\begin{example}
``Not mentioned'' means different things for different findings. For \textbf{Fracture} a silent report almost
always means no fracture (value \ca 0). For \textbf{Synovitis}, 15 of the 41 gold studies whose report was silent
were positive, so N gets a value of about \textbf{0.33}: ``we don't know, but a third of these knees have it''.
\end{example}

\subsection*{3. Read probabilities instead of text}

\begin{figure}[H]
\centering
\begin{tikzpicture}[x=1mm,y=1mm,
  seg/.style={draw=grid, fill=white, font=\scriptsize\ttfamily, minimum height=7mm, inner sep=1.5mm},
  ans/.style={draw=them, fill=whyfill, font=\scriptsize\ttfamily\bfseries, minimum height=7mm, inner sep=1.5mm}]
  \node[seg, anchor=west] (a) at (0,0) {\{"ACL": "};
  \node[ans, anchor=west] (b) at (a.east) {P};
  \node[seg, anchor=west] (c) at (b.east) {", "MCL": "};
  \node[ans, anchor=west] (d) at (c.east) {A};
  \node[seg, anchor=west] (e) at (d.east) {", "Medial Meniscus": "};
  \node[ans, anchor=west] (f) at (e.east) {A};
  \node[seg, anchor=west, draw=none] (g) at (f.east) {\ldots\ 12 findings + language};
  \node[seg, anchor=west, minimum height=4mm] (k1) at (0,-9) {};
  \node[font=\scriptsize, anchor=west] at (k1.east) {written by us};
  \node[ans, anchor=west, minimum height=4mm] (k2) at (30,-9) {};
  \node[font=\scriptsize, anchor=west] at (k2.east) {chosen by the model (we keep the probabilities of P, M, U, A, N)};
\end{tikzpicture}
\caption{We wrote the JSON answer ourselves and let the model fill in only the letters. At each gap we kept the
probabilities of all five letters (Figure~\ref{fig:nexttoken}). 13 tiny model calls per report instead of
\ca 110 generated tokens, and no parsing errors.}
\label{fig:prefill}
\end{figure}

The soft label is the probability-weighted value:
\[
\text{label} = \sum_{\text{letter}} p(\text{letter}) \cdot \text{value}(\text{letter})
\]

\begin{example}
Study K, ACL: $p$(P) = 0.95, $p$(M) = 0.02, $p$(U) = 0.02, $p$(A) = 0.01, $p$(N) = 0.00. \\
Label $= 0.95\cdot1 + 0.02\cdot0.2 + 0.02\cdot0.5 + 0.01\cdot0 = \mathbf{0.964}$. \\
Study K, Synovitis: $p$(N) = 0.97 $\Rightarrow$ label $\approx 0.97 \cdot 0.33 \approx \mathbf{0.32}$: ``unknown, prior
about one in three''.
\end{example}

\subsection*{4. A bigger model, made affordable}

\begin{figure}[H]
\centering
\begin{tikzpicture}[x=1mm,y=1mm]
  \fill[us!60] (0,0) rectangle (80.6,7);
  \node[font=\scriptsize, text=white] at (40,3.5) {shared instructions: 1,409 tokens (computed once)};
  \fill[them!70] (80.6,0) rectangle (100,7);
  \node[font=\scriptsize, text=white] at (90.3,3.5) {report};
  \draw[decorate, decoration={brace, amplitude=1.5mm}] (0,8) -- (100,8) node[midway, above=2mm, font=\scriptsize]
      {median prompt: 1,747 tokens};
  \node[font=\footnotesize, align=left, anchor=west] at (104,3.5) {only this part\\ is new per report};
\end{tikzpicture}
\caption{81\,\% of every prompt is the same instructions. Their internal state (the \term{KV cache}) was computed
once and reused for every report, which made the 32B model affordable on two T4s.}
\end{figure}

\paragraph{Safety rails.} Self-tests at start-up (the shortcut gives the same numbers as the plain model); the
58 gold reports first, and the 4,407-report run only if their macro AUC reached 0.84 (a \term{gate}); a time
limit of 10.5\,h with resume, because Kaggle kills sessions at 12\,h and a killed run saves nothing.

\section{Results}

\begin{figure}[H]
\centering
\begin{tikzpicture}
\begin{axis}[width=11cm, height=7.4cm, xmin=0.64, xmax=1.0,
    symbolic y coords={Fracture,Contusion,Baker's,Synovitis,Effusion,PF OA,Lat.~OA,Med.~OA,Lat.~meniscus,Med.~meniscus,MCL,ACL},
    ytick=data, y tick label style={font=\footnotesize}, xmajorgrids, xlabel={agreement with the 58 gold labels (AUC)},
    xlabel style={font=\footnotesize}, enlarge y limits=0.05, legend style={at={(0.98,0.02)}, anchor=south east, fill=white}]
  \addplot[grid, line width=1.8pt, forget plot] coordinates {(0.832,ACL) (0.99,ACL)};
  \addplot[grid, line width=1.8pt, forget plot] coordinates {(0.836,MCL) (0.989,MCL)};
  \addplot[grid, line width=1.8pt, forget plot] coordinates {(0.886,Med.~meniscus) (0.929,Med.~meniscus)};
  \addplot[grid, line width=1.8pt, forget plot] coordinates {(0.774,Lat.~meniscus) (0.872,Lat.~meniscus)};
  \addplot[grid, line width=1.8pt, forget plot] coordinates {(0.92,Med.~OA) (0.933,Med.~OA)};
  \addplot[grid, line width=1.8pt, forget plot] coordinates {(0.845,Lat.~OA) (0.839,Lat.~OA)};
  \addplot[grid, line width=1.8pt, forget plot] coordinates {(0.758,PF OA) (0.837,PF OA)};
  \addplot[grid, line width=1.8pt, forget plot] coordinates {(0.752,Effusion) (0.869,Effusion)};
  \addplot[grid, line width=1.8pt, forget plot] coordinates {(0.676,Synovitis) (0.774,Synovitis)};
  \addplot[grid, line width=1.8pt, forget plot] coordinates {(0.904,Baker's) (0.903,Baker's)};
  \addplot[grid, line width=1.8pt, forget plot] coordinates {(0.713,Contusion) (0.729,Contusion)};
  \addplot[grid, line width=1.8pt, forget plot] coordinates {(0.806,Fracture) (0.886,Fracture)};
  \addplot[only marks, mark=*, mark size=2.4pt, color=muted] coordinates {
    (0.832,ACL) (0.836,MCL) (0.886,Med.~meniscus) (0.774,Lat.~meniscus) (0.920,Med.~OA) (0.845,Lat.~OA)
    (0.758,PF OA) (0.752,Effusion) (0.676,Synovitis) (0.904,Baker's) (0.713,Contusion) (0.806,Fracture)};
  \addplot[only marks, mark=*, mark size=2.6pt, color=us] coordinates {
    (0.990,ACL) (0.989,MCL) (0.929,Med.~meniscus) (0.872,Lat.~meniscus) (0.933,Med.~OA) (0.839,Lat.~OA)
    (0.837,PF OA) (0.869,Effusion) (0.774,Synovitis) (0.903,Baker's) (0.729,Contusion) (0.886,Fracture)};
  \legend{{v3 labels (macro 0.809)}, {v4 labels (macro 0.879)}}
\end{axis}
\end{tikzpicture}
\caption{How well each label set agrees with the doctors, per finding. The threshold-aware v4 fixed ACL and MCL
(almost perfect: reports state ligament tears clearly), and helped effusion, fracture and PF OA. Synovitis and
contusion stayed weak.}
\label{fig:v3v4}
\end{figure}

\begin{innumbers}
Same image model, only the labels changed: \textbf{0.889 $\rightarrow$ 0.905} on the public leaderboard. On gold, the
model's ACL went 0.887 $\rightarrow$ 0.934, MCL 0.810 $\rightarrow$ 0.887, PF OA 0.799 $\rightarrow$ 0.900. The
biggest single gain of our whole project.
\end{innumbers}

\section{What the labeller actually answered}

\begin{figure}[H]
\centering
\begin{tikzpicture}
\begin{axis}[width=8.2cm, height=9cm, enlargelimits=false, colormap={bl}{color=(white) color=(us)},
    xtick={0,...,4}, xticklabels={P,M,U,A,N}, x tick label style={font=\footnotesize\bfseries},
    ytick={0,...,11}, yticklabels={ACL,MCL,Med.\ meniscus,Lat.\ meniscus,Med.\ OA,Lat.\ OA,PF OA,Effusion,Synovitis,Baker's,Contusion,Fracture},
    y tick label style={font=\footnotesize}, y dir=reverse, axis line style={draw=none}, tick style={draw=none},
    nodes near coords={\pgfmathprintnumber[fixed, precision=0]{\pgfplotspointmeta}}, nodes near coords style={font=\scriptsize, anchor=center},
    point meta min=0, point meta max=100, xticklabel pos=top]
  \addplot[matrix plot*, mesh/cols=5, point meta=explicit] table[x=x, y=y, meta=share] {figdata/letters.dat};
\end{axis}
\end{tikzpicture}
\caption{The labeller's most likely letter for each finding, in \% of the 4,349 report-only studies. Synovitis is
``not mentioned'' (N) in \synSilent\,\% of reports. MCL is P in only 1.4\,\% (the threshold wording was too strict:
6 of the 9 gold-positive MCL studies were answered M).}
\label{fig:letterheat}
\end{figure}

\begin{figure}[H]
\centering
\begin{tikzpicture}
\begin{axis}[ybar interval, width=10cm, height=4.6cm, xmin=0, xmax=1, ymin=0, xlabel={soft label for Effusion},
    ylabel={studies}, xlabel style={font=\footnotesize}, ylabel style={font=\footnotesize}, ymajorgrids,
    xtick={0,0.2,...,1}, xticklabel style={/pgf/number format/fixed}]
  \addplot[fill=us!70, draw=white] table[x=left, y=n] {figdata/effusion_hist.dat};
\end{axis}
\end{tikzpicture}
\caption{What the v4 labels look like for Effusion. Three peaks: clearly no (0, 38\,\%), clearly yes (1, 22\,\%), and
a large group at \textbf{0.44}: reports that mention an effusion without saying how much (answer U). The value of
U for Effusion, measured on the gold studies, is 0.44: a bit less than half of such knees had a moderate or large
effusion, so 0.44 is the honest label. Soft-label training
(Figure~\ref{fig:bce}) uses these in-between values directly.}
\end{figure}

\section{What reports cannot do}

\begin{keyidea}
A report-based label can only be as good as the report. Where reports are explicit (ACL, MCL) our labels reached
0.99. Where they are silent or vague (Synovitis, Contusion, Lateral OA) no prompt could fix it: \textbf{the
information is in the images, not in the text.} This is exactly where the winners' approach (reading the images
as well) made the difference.
\end{keyidea}

\begin{checkpoint}
\begin{itemize}
\done I can explain the five letters and how a soft label is computed from them.
\done I know why reading probabilities beats parsing generated text.
\done I know which findings reports label well and which they cannot.
\quiz Q7.1: Why is ``not mentioned'' not the same as ``absent''? Give a finding where the difference is large.
\quiz Q7.2: v4 labels improved ACL from 0.83 to 0.99 agreement. Why did the image model's gold ACL only go from
      0.89 to 0.93?
\quiz Q7.3: What does the ``gate'' protect against?
\end{itemize}
\end{checkpoint}
